\documentclass[11pt]{article}
\usepackage[T1]{fontenc}
\usepackage{lmodern,amsmath,amssymb,amsthm,booktabs,geometry,longtable,array,float,microtype,xcolor,hyperref}
\geometry{margin=1.05in}
\hypersetup{colorlinks=true,urlcolor=blue!60!black,linkcolor=blue!50!black,citecolor=blue!50!black,pdftitle={Rank-seven Grothendieck rings: a certified partial classification through multiplicity four},pdfauthor={}}
\newtheorem{theorem}{Theorem}
\newtheorem{proposition}[theorem]{Proposition}
\newtheorem{remark}[theorem]{Remark}
\newcommand{\C}{\mathbb C}
\newcommand{\Z}{\mathbb Z}
\newcommand{\FPdim}{\operatorname{FPdim}}
\title{Rank-seven Grothendieck rings: certified exclusions and models through multiplicity four}
\author{}
\date{}
\begin{document}
\maketitle
\begin{abstract}
We apply general categorification obstructions to the pinned, complete
fusion-ring census of rank seven and multiplicity at most four. Every
asserted exclusion has a replayable exact certificate; every asserted
realization has an explicit unitary fusion-category model and an exact
full-tensor identification. We include a rank-seven extension of the
formal-codegree test to noncommutative rings, a global-dimension-fifteen
obstruction, and complete oriented signed-sphericalization certificates
for 246 rings through multiplicity four. The remaining tensors are recorded as
\emph{undecided}. Three additional published AnyonWiki associators have
independent exact pentagon, invertibility, unit and rigidity checks.
Thus multiplicity one is complete, with thirteen realized rings;
the first undecided multiplicity is two.
\end{abstract}

\section{Scope and statements}
Let $R$ have distinguished basis $(b_0=1,b_1,\ldots,b_6)$ and multiplication
$b_i b_j=\sum_k N_{ij}^k b_k$. A \emph{Grothendieck ring} here means the
based ring of a complex fusion category; no pivotal, spherical, braided,
or unitary hypothesis is placed on a hypothetical categorification.
Its multiplicity is $\max_{i,j,k}N_{ij}^k$. The upstream
\emph{Fusion Ring Census} \cite{Census} provides the complete rank-seven
fusion-ring list through multiplicity four. We pin its source bytes,
source commit, row ordering, and their digest in the companion package.
This is a categorification study of that census; the list of undecided
tensors is not counted as a list of Grothendieck rings.
Multiplicities at least five lie outside the pinned source and the
claims of this manuscript.

\begin{theorem}[Certified rank-seven status through multiplicity four]
\label{thm:status}
The numbers of excluded, explicitly unitarily realized, and undecided
rank-seven fusion rings are given by Table~\ref{tab:status}.
The negative conclusions hold for \emph{every} complex fusion category.
The nineteen positive rows have the models in Table~\ref{tab:models},
matched against their complete fusion tensors. No conclusion is drawn
for the undecided rows.
\end{theorem}
\begin{table}[H]
\centering
\caption{Certified status of complete rank-seven fusion-ring census through exact multiplicity four. Undecided entries are not assumed categorifiable.}\label{tab:status}
\begin{tabular}{@{}r r r r r@{}}
\toprule
Exact multiplicity & Fusion rings & Excluded & Realized & Undecided\\
\midrule
1 & 43 & 30 & 10 & 3\\
2 & 319 & 291 & 5 & 23\\
3 & 1059 & 1035 & 1 & 23\\
4 & 2780 & 2751 & 0 & 29\\
\bottomrule
\end{tabular}
\end{table}

\begin{table}[H]
\centering
\caption{Sequential exclusion counts, charged only to the first successful criterion. Costs generally rise toward the bottom.}\label{tab:stages}
\begin{tabular}{@{}l r r r r@{}}
\toprule
Criterion & $m=1$ & $m=2$ & $m=3$ & $m=4$\\
\midrule
Near-group theorem & 2 & 1 & 1 & 1\\
Global dimension 15 & 1 & 0 & 0 & 0\\
Zero spectrum & 8 & 30 & 51 & 63\\
Closed rank-four/five subring & 0 & 25 & 125 & 235\\
Formal-codegree $d$-number & 15 & 179 & 788 & 2314\\
Subcategory dimension quotient & 0 & 5 & 9 & 8\\
Oriented dimension order & 4 & 51 & 61 & 130\\
\bottomrule
\end{tabular}
\end{table}

\begin{table}[H]
\centering
\caption{Unitary model constructions identified by exact full-tensor based-ring isomorphism. Each row denotes a ring, not a count of categories.}\label{tab:models}
\begin{tabular}{@{}l l@{}}
\toprule
Census ID & Explicit model\\
\midrule
\texttt{R7\_1\_1} & $\operatorname{Rep}(D_{16})$\\
\texttt{R7\_1\_6} & $\mathrm{SU}(2)_6$\\
\texttt{R7\_1\_9} & $\mathrm{PSU}(2)_{12}$\\
\texttt{R7\_1\_12} & $\mathrm{SO}(7)_2$\\
\texttt{R7\_1\_16} & $\operatorname{Rep}(D_{22})$\\
\texttt{R7\_1\_17} & $\mathrm{PSU}(2)_{13}$\\
\texttt{R7\_1\_24} & $\operatorname{Zest}(\mathrm{SU}(2)_6)$\\
\texttt{R7\_1\_25} & $\operatorname{Rep}(SD_{16})$\\
\texttt{R7\_1\_39} & $\operatorname{TY}(C_6)$\\
\texttt{R7\_1\_43} & $\operatorname{Vec}(C_7)$\\
\texttt{R7\_2\_24} & $\operatorname{Rep}(S_5)$\\
\texttt{R7\_2\_233} & $\operatorname{Rep}(H_{24})$\\
\texttt{R7\_2\_234} & $\mathcal C(A_5,1,A_4,1)$\\
\texttt{R7\_2\_297} & $\operatorname{Rep}(SL_2(\mathbb F_3))$\\
\texttt{R7\_2\_319} & $\operatorname{Rep}(C_{13}\rtimes C_3)$\\
\texttt{R7\_3\_1059} & $\operatorname{Rep}(C_{11}\rtimes C_5)$\\
\bottomrule
\end{tabular}
\end{table}


Multiplicity one is complete: its former residuals $R7\_1\_19$,
$R7\_1\_28$, and $R7\_1\_36$ have the AnyonWiki models described below.
The first undecided multiplicity is two, beginning with $R7\_2\_1$.
The exhaustive source-position convention is $R7\_m\_j$, where $j$
is the position \emph{among rows of exact multiplicity $m$} in the
pinned census file. This convention makes the positive-model
permutations and negative witnesses unambiguous.

\section{Necessary conditions and their order}
The checks in Table~\ref{tab:stages} run from top to bottom on surviving
tensors, and a tensor is charged to its first successful condition.
Near-group and small-dimension checks read a few coefficients; the
zero-spectrum and closed-subring passes are finite tensor scans;
exact Casimir factoring and subcategory-dimension arithmetic follow;
the complete oriented search is last. This order records precisely when
a criterion first contributes rather than hiding it in an unordered
union of tests. Multiplicity two is still undecided, so these stages
do \emph{not} establish a multiplicity threshold beyond which no new
criteria would be needed.
The oriented criterion first contributes at multiplicity one, so it
cannot be deferred to a higher multiplicity without losing a proof.

If $G$ is the group of invertible basis elements and the only other
object $X$ satisfies $X^2=\sum_{g\in G}g+\ell X$, then for $|G|=6$:
at $\ell=0$ the Tambara--Yamagami construction forces $G$ abelian
\cite{TY}; at $0<\ell<5$ in the commutative case,
Schopieray's near-group obstruction rules out categorification
\cite[Proposition~3.2]{Schopieray}. These statements apply to a
categorification without a pivotal assumption. The certificate records
the six group elements, the remaining object and $\ell$, and rechecks
every fusion-rule hypothesis.

The global-dimension-fifteen branch has six dimension-one simples and
one dimension-three simple: direct substitution into every product
establishes that $(1,\ldots,1,3)$ is its strictly positive dimension
character. Hence its Frobenius--Perron dimension is $6+3^2=15$.
The $pq$ classification \cite[Theorems~3.1 and~5.1]{EGO} makes a
complex fusion category of dimension $3\cdot5$ pointed in this case,
contradicting the dimension-three simple.

The zero-spectrum obstruction is the exact necessary criterion of
Liu--Palcoux--Ren \cite{LPR}; each certificate stores its nine basis
indices and the verifier rechecks all required zero/nonzero fusion
coefficients. A proper fusion-closed subset of basis elements would
generate a fusion subcategory. If its induced rank-four or rank-five
ring is excluded by the complete lower-rank classifications
\cite{Rank4,Rank5}, the ambient ring is excluded. The imported
model registries contain all fifteen rank-four and all seventeen
rank-five realized based-ring types, checked up to multiplicities
$128$ and $20$, respectively; all encountered subrings fall within
these ranges. The complete rank-six multiplicity-two classification
\cite{Rank6} supplies context and methodology, without being used as
an unstated rank-seven classification theorem.

\begin{proposition}[The formal-codegree test at rank seven]
\label{prop:dnumber}
Suppose that $R$ is a rank-seven fusion ring, possibly noncommutative.
Let $C=\sum_{i=0}^6 N_iN_i^{\mathsf T}$ be its regular Casimir matrix.
If an irreducible factor of the characteristic polynomial of $C$ is
not the minimal polynomial of a $d$-number, then $R$ cannot be the
Grothendieck ring of a complex fusion category.
\end{proposition}
\begin{proof}
In the commutative case, the eigenvalues of $C$ are the formal
codegrees of one-dimensional characters and must be $d$-numbers by
Ostrik \cite{OstrikCodegrees}. In the noncommutative case,
semisimplicity, $\dim_\C(R\otimes\C)=7$, and the positive
one-dimensional character imply the unique possible decomposition
$R\otimes\C\simeq\C^3\oplus M_2(\C)$.
There is only one two-dimensional irreducible module; it is fixed by
every field automorphism, so its formal codegree is a rational
algebraic integer and hence an integer. The Casimir acts on this
matrix block by twice that integer. Every remaining nonrational
Casimir eigenvalue comes from a one-dimensional character and is a
formal codegree, so Ostrik's condition applies. The implementation
factors the exact integral characteristic polynomial and tests the
coefficient-divisibility characterization of $d$-numbers. The replay
checker recomputes and factors that polynomial.
\end{proof}

If $S$ is a proper commutative fusion subring of a hypothetical
categorification, the quotient $\FPdim(R)/\FPdim(S)$ must be an
algebraic integer \cite{ENO}. The subcategory-divisibility stage
computes this quotient in the exact Frobenius--Perron number field,
records its characteristic polynomial, and certifies a nonintegral
coefficient. Its certificate is rederived from the complete tensor.

\section{Complete oriented tests at rank seven}
The last row of Table~\ref{tab:stages} uses the general, \emph{oriented}
signed-sphericalization construction established in \cite{Rank6}.
For an arbitrary hypothetical category $\mathcal C$, its canonical
sphericalization has two lifts of each simple object, interchanged
by a sign object $\varepsilon$. Taking the quotient by
$(\varepsilon+1)$ produces integer coefficients $T_{ij}^k$ with
$|T_{ij}^k|\le N_{ij}^k$ and
$T_{ij}^k\equiv N_{ij}^k\pmod 2$.
Unit, duality, cyclic Frobenius reciprocity, and ordered associativity
further constrain the quotient. The self-dual positive-indicator
normalization imposes the established square-channel congruences
modulo four. These conditions are necessary and cover every
categorification, whether pivotal or not; a signed solution is not
itself a categorification.

The rank-six proof of the general construction and the
dimension-order obstruction \cite{Rank6} are rank-independent. Here
their enumeration engine is extended to seven basis elements. It
enumerates the independent cyclic/dual-reflection coefficient orbits,
prunes only exact associativity equations with checked integer bounds,
and constructs \emph{every} signed quotient. For each quotient the
checker verifies either its adjoint parity, its gauge equivalence
with the original ring, or the failure of every admissible spherical
character by codegree and dimension-order tests. In particular the
dimension-order test requires the central-induction dimension
$\delta/f_\chi$ to lie in the $\Z$-span of the signed simple
dimensions, as required by \cite{OstrikPivotal,Rank6}.

\begin{proposition}[Verified exhaustive branches]
The oriented-dimension-order criterion excludes four, fifty-one,
sixty-one, and one hundred thirty previously surviving rank-seven
rings of exact multiplicities one through four, respectively.
For $R7\_1\_5$, $R7\_1\_23$, $R7\_1\_33$, and $R7\_1\_38$,
the complete signed enumerations have respectively $2,4,4,8$
associative quotient tables. Every table in each of these 246
certified exhaustive searches is eliminated by necessary conditions
for arbitrary complex fusion categories.
\end{proposition}
\begin{proof}
The machine-readable proof tree for each of the 246 rings gives its full
dimension-character factorization, every obstruction, the ordered
signed-orbit domain, and one checked branch for every enumerated
table. The independent verifier rebuilds the domains from the source
tensor, repeats the bounded exhaustive enumeration, checks the
\emph{entire} list of surviving tables for equality to the recorded
list, and recomputes each branch witness. The search used to discover
these finite trees is logically separate from their exhaustive replay;
timed-out or unfinished searches generate no negative decision.
The enumeration returns the four displayed multiplicity-one branch
counts and no uncovered branch for any of the 246 certificates.
The rank-independent necessity of the search is the
sphericalization theorem cited above.
\end{proof}

For the three surviving multiplicity-one rows the current general
filters provide no exclusion. Publicly proposed numerical categorical
data \cite{Vercleyen,VercleyenInvariants} provide useful leads; no
unverified numerical solution is promoted here to a constructive
model. This same standard applies to the larger residual sets at
multiplicities two, three, and four.

\section{Exact constructions and replay}
The three formerly undecided multiplicity-one rings match AnyonWiki
$\mathrm{FR}^{7,2}_3$, $\mathrm{FR}^{7,2}_{12}$, and
$\mathrm{FR}^{7,4}_3$, respectively, by full based-ring isomorphisms.
For each, we use the published category with code $(7,1,d,i,1,0,1)$
in the Vercleyen--Slingerland census \cite{AnyonWiki2026}.
The symbolic associators supplied through TensorCategories.jl lie over
$\mathbb Q(\zeta_8)$, $\mathbb Q(\zeta_{28})$, and
$\mathbb Q(\zeta_{16})$, respectively. Independent rational-polynomial
replay checks all $9{,}856$, $45{,}061$, and $9{,}856$ admissible
pentagon equations, all $585$, $837$, and $585$ nonempty associator
blocks for invertibility, and the unit and rigidity constraints.
These checks establish complex fusion categories directly, with no
pivotality assumption. The AnyonWiki category list additionally records
unitary gauges. Only existence of a model is used: neither completeness
nor uniqueness of the external category census is required.
The supplementary evidence records the exact category codes,
cyclotomic coefficients, source hashes and basis permutations.

The positive model catalogue constructs complete fusion tensors from
the cyclic group, ordinary finite-group character products, the
Tambara--Yamagami rule, $SU(2)_k$ truncation, its even parts, odd
metaplectic rules, and the categorical $\Z/2$ zest of $SU(2)_6$
\cite{TY,KO,Metaplectic,Zest,Vercleyen}. Character products for
$\operatorname{Rep}(S_5)$, $\operatorname{Rep}(SL_2(\mathbb F_3))$,
$\operatorname{Rep}(C_{13}\rtimes C_3)$ and
$\operatorname{Rep}(C_{11}\rtimes C_5)$ are computed with exact
characters in the bundled software. Two further complete model tensors
come from the published noncommutative integral appendix of
\cite[Theorem~1.6 and Section~6.2]{Integral}: the representation
category of its dimension-24 Kac algebra and the group-theoretical
category $\mathcal C(A_5,1,A_4,1)$. The published full tensors and
their byte digest are pinned in the package. For each constructed or
published model, the program validates the fusion axioms, then checks
all $6!$ unit-preserving basis permutations against every compatible
census tensor. A positive record carries its explicit permutation.
Different categorical structures may share a based ring; the count
is of rings, not categories.

The complete package contains the pinned gzip source, analysis code,
negative proof trees, model tensors, the stage-count summary and a
second entry point, \texttt{verify.py}. Run
\texttt{python analyze.py --max-m 4 --dnumber} and then
\texttt{python verify.py}. The verifier rechecks the source digest and
every row, every full-tensor model isomorphism, each stage total and
each exclusion certificate, including the exhaustive signed lists.
Running \texttt{python -O verify.py} gives the same checks without
relying on Python's optional \texttt{assert} statements.

\medskip\noindent\textbf{AI disclosure.} GPT-6 Astra assisted with
research, symbolic exploration, code development, and drafting.

\begin{thebibliography}{99}
\bibitem{AnyonWiki2026} G.~Vercleyen and J.~Slingerland,
\emph{On Low-Rank Multiplicity-Free Complex Fusion Categories},
\emph{Commun. Math. Phys.} \textbf{407}, 231 (2026),
DOI:10.1007/s00220-026-05705-2;
\url{https://anyonwiki.github.io/pages/Lists/losmffc.html}.
\bibitem{Census} S. Palcoux, \emph{Fusion Ring Census},
\url{https://github.com/sebastienpalcoux/fusion-ring-census},
commit \texttt{b86eb6dc6eeeafe5fc5f6906d7b074e0322e9726}.
\bibitem{Rank4} \emph{Rank-four Grothendieck rings through multiplicity 128}, manuscript and exact-certificate supplement (2026).
\bibitem{Rank5} \emph{Rank-five Grothendieck rings through multiplicity twenty}, manuscript and exact-certificate supplement (2026).
\bibitem{Rank6} \emph{Rank-six fusion rings: categorification criteria and models through multiplicity eight}, manuscript and exact-certificate supplement (2026).
\bibitem{Integral} M. A. Alekseyev, W. Bruns, J. Dong and S. Palcoux,
\emph{Classifying integral Grothendieck rings up to rank 5 and beyond},
accepted in \emph{Letters in Mathematical Physics}; arXiv:2507.07023.
\bibitem{EGO} P. Etingof, S. Gelaki and V. Ostrik,
\emph{Classification of fusion categories of dimension $pq$},
\emph{Int. Math. Res. Not.} \textbf{2004}, 3041--3056; arXiv:math/0304194.
\bibitem{ENO} P. Etingof, D. Nikshych and V. Ostrik,
\emph{On fusion categories}, \emph{Ann. of Math.} \textbf{162} (2005), 581--642; arXiv:math/0203060.
\bibitem{KO} A. Kirillov Jr. and V. Ostrik,
\emph{On a $q$-analogue of the McKay correspondence and the ADE classification of $\widehat{\mathfrak{sl}}_2$ conformal field theories},
\emph{Adv. Math.} \textbf{171} (2002), 183--227; arXiv:math/0101219.
\bibitem{LPR} Z. Liu, S. Palcoux and Y. Ren,
\emph{Classification of Grothendieck rings of complex fusion categories of multiplicity one up to rank six},
\emph{Lett. Math. Phys.} \textbf{112} (2022), article 54; arXiv:2010.10264.
\bibitem{Metaplectic} P. Bruillard, C. Galindo, S.-H. Ng, J. Y. Plavnik, E. C. Rowell and Z. Wang,
\emph{On the classification of weakly integral modular categories},
\emph{J. Pure Appl. Algebra} \textbf{220} (2016), 2364--2388; arXiv:1411.2313.
\bibitem{OstrikCodegrees} V. Ostrik,
\emph{On formal codegrees of fusion categories},
\emph{Math. Res. Lett.} \textbf{16} (2009), 895--901; arXiv:0810.3242.
\bibitem{OstrikPivotal} V. Ostrik,
\emph{Pivotal fusion categories of rank 3},
\emph{Mosc. Math. J.} \textbf{15} (2015), 373--396; arXiv:1309.4822.
\bibitem{Schopieray} A. Schopieray,
\emph{Categorification of integral group rings extended by one dimension},
\emph{J. Lond. Math. Soc.} \textbf{108} (2023); DOI:10.1112/jlms.12789; arXiv:2208.07319.
\bibitem{TY} D. Tambara and S. Yamagami,
\emph{Tensor categories with fusion rules of self-duality for finite abelian groups},
\emph{J. Algebra} \textbf{209} (1998), 692--707.
\bibitem{Zest} C. Delaney, C. Galindo, J. Y. Plavnik, E. C. Rowell and Q. Zhang,
\emph{Braided zesting and its applications},
\emph{Comm. Math. Phys.} \textbf{386} (2021), 1--55; arXiv:2005.05544.
\bibitem{Vercleyen} G. Vercleyen,
\emph{On low-rank multiplicity-free fusion categories}, thesis (2024),
arXiv:2405.20075.
\bibitem{VercleyenInvariants} G. Vercleyen,
\emph{Tables of practical invariants for distinguishing multiplicity-free fusion categories up to rank 7}, arXiv:2507.00652.
\end{thebibliography}
\end{document}
